% Shared exercise chunk: l1-rope-collisions
\begin{exercisechunk}
\item \textbf{RoPE collisions and windowed attention.}
\label[exercise]{ex:rope-collisions}
\exerciseprofile{2}{2}{3}{\exproof\quad\excalculation\quad\exconstruction}{%
Characterize periodic address collisions and prove that a minimal attention
window removes them.}
Use the construction and attention output from
\cref{ex:rope,ex:rope-retrieval}. Hold $d$ fixed, let the sequence length grow,
choose an integer period $P>d$, set $\omega=2\pi/P$, and initially use
$S_i=[i]$. A \emph{collision} occurs when another source has the same score as
the requested source $i-d$.
\begin{enumerate}[(a),beginpenalty=10000]
\item Determine the set of colliding source positions.
\item Show that the first collision occurs at $i=d+P+1$.
\item At this first collision, show that
\[
 \lim_{\lambda\to\infty}o_i(\lambda)
 =\frac{v_1+v_{P+1}}2.
\]
\item Now use the window
$S_i=\{j\in[i]:i-j<w\}$, where $w$ is a positive integer. Show that the
smallest width retaining $i-d$ is $w=d+1$. For this width, prove that $i-d$ is the unique
maximizer and that $\lim_{\lambda\to\infty}o_i(\lambda)=v_{i-d}$ for every
$i>d$, however long the sequence becomes.
\end{enumerate}
\end{exercisechunk}
